% TOP_PROBLEM: {"id": 12, "title": "Goldbach's Conjecture", "category_id": 1, "category": {"id": 1, "name": "number_theory", "display_name": "Number Theory", "description": "Properties of integers, prime numbers, Diophantine equations.", "slug": "number-theory", "order_index": 1, "created_at": "2026-07-31T15:26:25.670Z"}, "status": "open"}
% FIRST_POSTED: 2026-09-25
% ENTRY_KIND: statement_only
% !TeX program = lualatex
% Definition and sources only; this file is not an LLM solution attempt.
\documentclass[11pt]{article}
\usepackage[margin=25mm]{geometry}
\usepackage{fontspec}
\setmainfont{Latin Modern Roman}
\usepackage{amsmath,amssymb,mathtools}
\usepackage{unicode-math}
\setmathfont{Latin Modern Math}
\usepackage{xurl,hyperref}
\hypersetup{colorlinks=true,urlcolor=blue}
\setcounter{secnumdepth}{0}
\setlength{\parindent}{0pt}
\setlength{\parskip}{5pt}
\setlength{\emergencystretch}{3em}
\begin{document}
\section*{\texorpdfstring{Goldbach's Conjecture}{Goldbach's Conjecture}}\label{12}
\subsection{Definitions and mathematical statement}
With $\mathcal P$ denoting the positive primes, the statement is
\[
 \forall n\in2{\mathbb{Z}},\quad n>2\quad\Longrightarrow\quad
 \exists p,q\in\mathcal P\text{ with }n=p+q.
\]
The primes are not required to be distinct, so $4=2+2$ is allowed. Every even integer is quantified, not merely all sufficiently large integers or a density-one subset.
\par\smallskip\textit{Formulation and concept references:} \href{https://t5k.org/notes/conjectures/}{[S1]}.

\subsection{Short English statement}
Can every even integer larger than two be written as the sum of two primes?
\subsection{Sources}
\begin{itemize}
\item[S1]C. K. Caldwell, "Prime Conjectures and Open Questions," The PrimePages, https://t5k.org/notes/conjectures/.
\url{https://t5k.org/notes/conjectures/}.
\textit{Formulation source.}
\item[S2]The exceptional set of the Goldbach problem.
\url{https://arxiv.org/html/2607.27282v2}.
\textit{Further reference; not a proof endorsement.}
\item[S3]Theorem (1+1.9) on the Goldbach Conjecture.
\url{https://arxiv.org/abs/2606.05224}.
\textit{Further reference; not a proof endorsement.}
\item[S4]Restricted Goldbach Sums in Arithmetic Progressions: Local Obstructions, an Explicit Almost-All Framework, and a General-Modulus Extension.
\url{https://www.preprints.org/manuscript/202607.2077}.
\textit{Further reference; not a proof endorsement.}
\end{itemize}
\end{document}
